%% annex/p06/constructibility-forcing.tex --- proofs and sketches in body
%% order for the migrated constructibility/forcing chapter.

\begin{annexitem}{annex:p6.force.1}{thm:p06-constructible-class}{The constructible class}
\sketchnote{the level-by-level definability and limit-stage bookkeeping are
omitted; see \cite[Chapter~13]{Jech2003}.}
Transfinite induction gives that every \(L_\alpha\) is transitive.  The same
induction yields \(\alpha\in L_{\alpha+1}\): earlier ordinals and their
membership relation are available at preceding stages, and the next ordinal
is definable there.  Taking the union proves both claims for \(L\).
\end{annexitem}

\begin{annexitem}{annex:p6.force.2}{thm:p06-ctm-strength}{Strength of the countable-transitive-model hypothesis}
\sketchnote{the arithmetic-to-membership translation, absoluteness
calculation, and transport of the derivability conditions are omitted.}
The set-sized structure \((M,\in)\) witnesses metatheoretic consistency of
\(\mathsf{ZFC}\) by soundness.  Under the finite-ordinal interpretation, the
standard consistency sentence is true; arithmetic absoluteness gives
\((M,\in)\entails\mathsf{ZFC}+\mathsf{Con}(\mathsf{ZFC})^\in\).
Soundness yields consistency of that stronger theory.  Transported second
incompleteness says it cannot prove its own consistency, so the model-existence
hypothesis cannot follow from bare \(\mathsf{Con}(\mathsf{ZFC})\).
\end{annexitem}

\begin{annexitem}{annex:p6.force.3}{thm:p06-bounded-absoluteness}{Bounded absoluteness for the fixed transitive set}
\sketchnote{the Boolean-connective cases of the formula induction are omitted.}
Induct on the bounded formula.  Atomic membership and equality are absolute
because \(M\) is transitive.  For
\(\exists x\in a\,\psi(x,\vec a)\), every possible witness lies in \(M\)
when \(a\in M\), and the induction hypothesis transfers \(\psi\); the
universal bounded case is dual.
\end{annexitem}

\begin{annexitem}{annex:p6.force.4}{thm:p06-generic-filter-existence}{Existence of a generic filter under the fixed hypothesis}
\begin{proof}
Enumerate \(M\)'s dense sets as \(D_0,D_1,\ldots\).  Choose
\(p_0\leq p\) and \(p_{n+1}\leq p_n\) in \(D_n\).  Let
\[
  G=\{q\in\mathbb P:\exists n\,(p_n\leq q)\}.
\]
\(p\in G\), \(G\) is upward closed, \(p_n\) commonly strengthens
two members, and \(p_{n+1}\in G\cap D_n\).  Thus \(G\) is generic.
\end{proof}
\end{annexitem}

\begin{annexitem}{annex:p6.force.5}{thm:p06-check-name-evaluation}{Evaluation of check names}
\sketchnote{the rank bookkeeping placing each recursively formed name in
\(M\cap V^{\mathbb P}\) is omitted.}
Induct on the rank of \(x\).  A nonempty poset filter contains
\(1_{\mathbb P}\), so
\[
  (\check x)^G
  =\{(\check y)^G:y\in x\}
  =\{y:y\in x\}=x.
\]
Transitivity and Replacement inside \(M\) carry the simultaneous membership
claim through the same induction.
\end{annexitem}

\begin{annexitem}{annex:p6.force.6}{thm:p06-semantic-forcing-monotone}{Monotonicity of semantic forcing}
\begin{proof}
Let \(G\ni q\) be generic with a prescribed assignment.  Since \(q\leq p\),
upward closure gives \(p\in G\); the hypothesis makes \(\varphi\) true in
\(M[G]\).  Arbitrary \(G\) and assignment give the defining condition at \(q\).
\end{proof}
\end{annexitem}

\begin{annexitem}{annex:p6.force.7}{thm:p06-cohen-new-real}{One Cohen real}
\begin{proof}
Let
\[
  \dot c=\{\langle\check{\langle n,i\rangle},p\rangle:
    n\in\dom(p),\ p(n)=i\}.
\]
This name lies in \(M\), and
\[
  (\dot c)^G
  =\{\langle n,i\rangle:\exists p\in G\;p(n)=i\}
  =\bigcup G=c,
\]
so \(c\in M[G]\).  Directedness makes \(c\) single-valued: two conditions in
\(G\) assigning conflicting values have no common strengthening.  For each \(n\),
\(D_n=\{p:n\in\dom(p)\}\) belongs to \(M\) and is dense, so genericity makes
\(c\) total.  Given \(b\in M\cap2^\omega\), conditions disagreeing with \(b\)
form a dense member of \(M\): extend at a fresh coordinate with the opposite
value.  Meeting it yields \(c\neq b\).  Therefore \(c\in M[G]\setminus M\).
\end{proof}
\end{annexitem}

\begin{annexitem}{annex:p6.force.9}{thm:p06-ccc-preservation}{Cardinal and cofinality preservation by ccc forcing}
\begin{proof}
Work in \(M\), using Choice.  First note that if \(D\in M\) is dense below
some \(d\in G\), then
\[
  D\cup\{q\in\mathbb P:q\text{ is incompatible with }d\}
\]
is a globally dense member of \(M\).  Genericity and \(d\in G\) therefore
give \(G\cap D\neq\varnothing\).

If \(M[G]\) collapsed an uncountable
\(M\)-cardinal \(\kappa\), there would be an ordinal
\(\lambda<\kappa\), a surjection \(f:\lambda\to\kappa\) in \(M[G]\), a
name \(\nu_f\in M\) with \(\nu_f^G=f\), and \(p\in G\) such that
\(p\forces \nu_f:\check\lambda\twoheadrightarrow\check\kappa\).  For each
\(\alpha<\lambda\), choose below \(p\) a maximal antichain deciding
\(\nu_f(\check\alpha)\).  The ccc makes it countable; if \(X_\alpha\) is its set
of possible values, the downward closure of that antichain is dense below
\(p\).  The observation above shows that the interpreted map satisfies
\[
  \operatorname{ran}(f)\subseteq
  \bigcup_{\alpha<\lambda}X_\alpha.
\]
In \(M\), the union has size at
most \(\max((\abs\lambda)^M,\aleph_0)<\kappa\), impossible.  The finite cardinals
and \(\omega\) are absolute because the extension has the same ordinals.

Now let \(\mu=\operatorname{cf}^{M}(\delta)\).  A ground-model cofinal map
of type \(\mu\) remains cofinal because the ordinals and their order are
unchanged, so cofinality cannot increase.  If a cofinal map in \(M[G]\) had
infinite domain ordinal \(\lambda<\mu\), choose a name and a condition in
\(G\) forcing this, and repeat the deciding-antichain construction below that
condition.  Since \(\mu\) is regular, the resulting union of possible values
has size \(\max((\abs\lambda)^M,\aleph_0)<\mu\) in \(M\), hence is bounded in
\(\delta\), contradicting cofinality.  A finite range is bounded in every
limit ordinal, so ground cofinality \(\omega\) cannot drop to a finite value;
successor and zero cofinalities are fixed by the unchanged ordinals.
\end{proof}
\end{annexitem}

\begin{annexitem}{annex:p6.force.10}{thm:p06-nice-name-counting}{Nice-name reduction and counting}
\begin{proof}
Fix \(p\forces\nu\subseteq\check\omega\).  For each \(n<\omega\), choose
below \(p\) a maximal antichain \(A_n\) among conditions forcing
\(\check n\in\nu\).  Put
\[
  \zeta=\bigcup_{n<\omega}
  \{\langle\check n,a\rangle:a\in A_n\}.
\]
For any generic \(H\ni p\), a member of \(A_n\cap H\) forces
\(n\in\zeta^H\cap\nu^H\).  Conversely, if \(n\in\nu^H\), the truth lemma
gives \(q\in H\) forcing \(\check n\in\nu\).  Directedness supplies
\(d\in H\) with \(d\leq p,q\), which still forces membership.  Maximality
makes the common extensions of \(d\) with members of \(A_n\) dense below
\(d\).  Adjoining every condition incompatible with \(d\) makes this set a
globally dense member of \(M\).  Since \(d\in H\), genericity cannot meet the
incompatible part, so it supplies a common extension in \(H\); upward closure
then puts the corresponding member of \(A_n\) in \(H\).  Hence
\(n\in\zeta^H\).  Thus
\(p\) semantically forces \(\nu=\zeta\), and the forcing theorem yields
\(p\forces\nu=\zeta\).

Under ccc each \(A_n\) is countable.  There are at most
\(\lambda=(\abs{\mathbb P}^{\aleph_0})^M\) countable subsets of
\(\mathbb P\) in \(M\), and a countable sequence of them has the same bound.
For any \(x\subseteq\omega\) in \(M[G]\), choose \(\nu\in M\) with
\(x=\nu^G\).  The truth lemma supplies \(p\in G\) forcing
\(\nu\subseteq\check\omega\), so the local result gives a nice \(\zeta\)
with \(x=\zeta^G\).  Thus the nice names cover all reals; ccc preservation
keeps \(\lambda\) a cardinal, proving the bound.
\end{proof}
\end{annexitem}

\begin{annexitem}{annex:p6.force.11}{thm:p06-many-cohen-continuum}{Forcing the size of the continuum}
\begin{proof}
For \(\alpha<\kappa\), let
\(c_\alpha(n)=(\bigcup G)(\alpha,n)\).  Coordinate-density makes each
\(c_\alpha\) total.  For distinct \(\alpha,\beta\), conditions assigning
opposite values at a fresh common coordinate are dense, so
\(c_\alpha\neq c_\beta\); hence the
extension has at least \(\kappa\) reals.  The poset has cardinality \(\kappa\),
is ccc by Theorem~\ref{thm:p06-add-ccc}, and preserves \(\kappa\) by
Theorem~\ref{thm:p06-ccc-preservation}.  Nice-name counting and
\(\kappa^{\aleph_0}=\kappa\) give at most \(\kappa\) reals.  Therefore the
continuum is \(\kappa\); for ground \(\mathsf{GCH}\) and
\(\kappa=\aleph_2\), CH fails.
\end{proof}
\end{annexitem}

\begin{annexitem}{annex:p6.force.12}{thm:p06-relativization-derivations}{Relativization of derivations}
\begin{proof}
Strengthen the claim to every subderivation with open assumptions \(\Gamma\),
conclusion \(\chi\), and any finite variable list containing
\(\operatorname{FV}(\Gamma\cup\{\chi\})\).  Induct on the derivation under
the list's universal quantifiers and \(L\)-membership guards.  If an induction
hypothesis contains a variable absent from the target node, instantiate it at
the \(\mathsf{ZF}\)-provable witness \(\varnothing\in L\), removing the
unnecessary quantifier and guard.

Relativization fixes atoms and \(\bottom\) and commutes with propositional
connectives.  Thus assumption, propositional, and classical steps reproduce
the same rules, with exactly the originally marked nonvacuous discharges.
Equality introduction is unchanged; variable substitution commutes with
relativization up to alpha-equivalence, giving equality elimination.  The
explicit alpha rule remains available because the inserted guard is
parameter-free.

For \(\forall E\), instantiate
\(\forall x(x\in L\to\chi^L)\) and use the term-variable's guard.  For
\(\exists I\), conjoin that guard before introduction.  A \(\forall I\)
step temporarily assumes \(x\in L\), derives the relativized premise, then
discharges that occurring assumption before generalization; when \(x\) is
absent from the premise, first conjoin and project so that discharge remains
nonvacuous.  An \(\exists E\) step opens the relativized existential, projects
both \(y\in L\) and the minor assumption from their conjunction, and applies
the induction hypothesis.  The original eigenvariable remains absent from the
conclusion, all other open assumptions, and the existential premise because
relativization adds no free variables.  Capture-avoiding quantifier
substitution may rename bound variables; the explicit alpha rule connects the
resulting variants.  Finally discharge the relativized open assumptions and
introduce the listed universal quantifiers.
\end{proof}
\end{annexitem}

\begin{annexitem}{annex:p6.force.13}{thm:p06-ch-relative-consistency}{Positive relative consistency of CH}
\begin{proof}
If axioms \(\sigma_1,\ldots,\sigma_k\) of
\(\mathsf{ZFC}+\mathsf{GCH}+(V=L)\) derived \(\bottom\),
Theorem~\ref{thm:p06-relativization-derivations} would give
\[
  \mathsf{ZF}\derives_{\mathsf{ND}_{\mathsf{FOL}}}
  (\sigma_1^L\land\cdots\land\sigma_k^L)\to\bottom.
\]
The cited constructibility scheme supplies each \(\sigma_i^L\), contradicting
consistency of \(\mathsf{ZF}\).

Consistency of \(\mathsf{ZFC}\) implies consistency of \(\mathsf{ZF}\), so
the constructible theory is consistent and proves \(\mathsf{CH}\) as the
\(\alpha=0\) instance of \(\mathsf{GCH}\).  If
\(\mathsf{ZFC}+\mathsf{CH}\) were inconsistent, its derivation must have an
open \(\mathsf{CH}\) occurrence, since otherwise \(\mathsf{ZFC}\) itself
would be inconsistent.  Nonvacuous negation introduction would give
\(\mathsf{ZFC}\derives_{\mathsf{ND}_{\mathsf{FOL}}}\neg\mathsf{CH}\);
monotonicity would make the
consistent constructible theory prove both sides.
\end{proof}
\end{annexitem}

\begin{annexitem}{annex:p6.force.14}{thm:p06-arithmetized-ch-consistency}{Arithmetized consistency implications}
\sketchnote{the coding of both derivation transformations, including the
primitive-recursive translation between the fixed calculus and the calculus
used by the Boolean-valued source, is omitted.}
Relativization and the Boolean-valued consistency discharge transform any
target-theory contradiction into a
\(\mathsf{ZFC}\)-contradiction.  Their syntax recursions and finite derivation
operations are primitive recursive.  Representing those transformations under
the fixed proof predicate yields both implications in \(\mathsf{PA}\).
\end{annexitem}

\begin{annexitem}{annex:p6.force.15}{thm:p06-ch-nonderivability}{Nonderivability of CH and its negation}
\begin{proof}
If \(\mathsf{ZFC}\derives_{\mathsf{ND}_{\mathsf{FOL}}}\mathsf{CH}\),
monotonicity and the assumption rule
make \(\mathsf{ZFC}+\neg\mathsf{CH}\) derive both sides, hence
\(\bottom\); the cited negative relative-consistency theorem contradicts the
hypothesis.  If
\(\mathsf{ZFC}\derives_{\mathsf{ND}_{\mathsf{FOL}}}\neg\mathsf{CH}\),
the same argument
makes \(\mathsf{ZFC}+\mathsf{CH}\) inconsistent, contradicting
Theorem~\ref{thm:p06-ch-relative-consistency}.
\end{proof}
\end{annexitem}
